\begin{textsample}{POS Dim 2 – 2017\_c5 – Score -1.00 – 2017\_c5/enade\_2017\_tec\_gest\_prod\_99ff610856.txt}  \label{ex:f2_pos_387}
CAPÍTULO 5 DISTRIBUIÇÃO DOS CONCEITOS 5.1 PANORAMA NACIONAL DA DISTRIBUIÇÃO DOS CONCEITOS O Conceito Enade é um indicador de qualidade que avalia os cursos por intermédio dos desempenhos dos estudantes a partir dos resultados do Enade.
Os valores possíveis vão de 1 ( pior situação ) a 5 ( melhor situação ).
Os cursos com apenas um ou sem nenhum concluinte participante não obtêm o Conceito Enade, ficando Sem Conceito ( SC ).
Na Tabela 5.1, são apresentadas a quantidade e a distribuição de cursos de Tecnologia em Gestão da Produção Industrial participantes do Enade/2017, por faixa de conceito e Grande Região.
Enfatiza -se, mais uma vez, que a diferença entre os cursos tabulados neste capítulo e no capítulo 2 corresponde aos cursos Sem Conceito ( SC ), ou seja àqueles, em princípio, sem alunos concluintes que houvessem participado da prova ou àqueles que tivessem tão somente um aluno concluinte¹⁸.
Observando -se os dados da Tabela 5.1, nota -se que, dos 62 cursos participantes, 20 ( 32,3 \% ) classificaram -se com conceito 4, o valor modal.
Esse foi também o conceito modal na maioria das regiões : Sudeste ( 29,4 \% ), Sul ( 47,4 \% ), e Centro-Oeste ( 100 \% ).
Na região Norte, a moda foi o conceito 3, atribuído para 50,0 \% dos cursos da região, e na região Nordeste, o conceito modal foi 2 ( 75,0 \% ).
O conceito 3 foi o segundo mais frequente em nível nacional ( 27,4 \%, correspondendo a 17 cursos ), e o conceito 2, o terceiro ( 16,1 \%, correspondendo a dez cursos ).
Houve, ainda, dois cursos ( 3,2 \% ) que receberam conceito 5 e mais dois cursos ( 3,2 \% ) que receberam conceito 1.
Dos 62 cursos de Tecnologia em Gestão da Produção Industrial, 11 ( 17,7 \% ) ficaram Sem Conceito ( SC ). ¹⁸ Estes cursos com somente um concluinte não foram considerados no capítulo 2.
A região Norte participou com quatro cursos ou 6,5 \% do total nacional.
Desses, dois cursos receberam o conceito 3, o que equivale a 50,0 \% do total regional, além de ser o conceito modal na região, como já comentado.
A um curso ( 25,0 \% ) atribuiu -se conceito 2 e um curso ficou Sem Conceito ( SC ), nenhum curso recebeu conceitos 1, 4 ou 5 nessa região.
A região Nordeste também participou com quatro cursos ou 6,5 \% do total nacional.
Nessa região, três cursos ( 75,0 \% em termos regionais ) obtiveram conceito 2, o conceito modal para a região, como já comentado.
Os demais conceitos não foram atribuídos a nenhum curso.
Ainda nessa região, um curso ficou Sem Conceito ( SC ).
Dos 34 cursos participantes da região Sudeste ( 54,8 \% do total nacional ), o conceito 4 foi o conceito modal, como já anteriormente assinalado, tendo sido obtido por dez cursos ( 29,4 \% ).
O conceito 1 foi atribuído a dois cursos ( 5,9 \% ), e o conceito 2, a quatro cursos ( 11,8 \% ).
Receberam os conceitos 3 e 5, respectivamente, nove cursos ( 26,5 \% ) e um curso ( 2,9 \% ).
Nessa região, oito cursos ( 23,5 \% ) ficaram Sem Conceito ( SC ).
Os 19 cursos da região Sul corresponderam a 30,6 \% do total nacional.
Como já indicado, houve predominância do conceito 4, atribuído a 47,4 \% dos cursos da região ( conceito modal ), ou seja, atribuído a nove dos 19 cursos participantes na região Sul.
O conceito 2 foi atribuído a dois cursos ( 10,5 \% ), o conceito 3, a seis cursos ( 31,6 \% ) e o conceito 5, a um curso ( 5,3 \% ).
Um curso ficou Sem Conceito ( SC ) e nenhum recebeu conceito 1 na região Sul.
Já o único curso participante na região Centro-Oeste ( 1,6 \% do total nacional ), recebeu conceito 4.
Nessa região, nenhum curso ficou Sem Conceito ( SC ) ou recebeu os demais conceitos atribuídos.
No Gráfico 5.1, é apresentada a distribuição cumulativa do Conceito Enade, segundo a Grande Região.
Em um gráfico de distribuição cumulativa, em cada valor de abscissa ( eixo x ), o valor apresentado no eixo das ordenadas ( eixo y ) é o valor acumulado da variável.
Por exemplo, a região Nordeste ( linha verde ) apresenta 100 \% dos cursos com conceito 2, já que a sua poligonal cumulativa parte de 0 \% no conceito 1 e alcança 100 \% no conceito 2.
A região Sul ( linha roxa ), por outro lado, apresenta um pouco mais de 11 \% dos cursos na mesma situação : com conceito 2 ou abaixo.
Nas regiões Norte ( linha azul ), todos os cursos têm conceitos 3 ou abaixo disso, já que a poligonal cumulativa alcança 100 \% para o conceito 4.
Em linhas gerais, em um gráfico de distribuição cumulativa, poligonais mais à esquerda representam Grandes Regiões com pior distribuição, e poligonais mais à direita, Grandes Regiões com melhores conceitos.
A região Nordeste ( linha verde ) apresenta os cursos com conceitos mais baixos, ao passo que a região Sul ( linha roxa ) apresenta as distribuições com valores maiores.
Cursos em IES situadas na região Centro-Oeste ( linha preta ) apresentam uma situação peculiar, receberam apenas conceito 4, mas mesmo tendo a poligonal mais à direita que a região Sul ( linha roxa ), esta supera aquela por ter uma baixa porcentagem de valores de conceitos baixos e apresenta cursos com conceito máximo. 5.2 CONCEITOS POR CATEGORIA ADMINISTRATIVA, POR MODALIDADE DE ENSINO E GRANDE REGIÃO Na Tabela 5.2, é apresentada a distribuição dos cursos de Tecnologia em Gestão da Produção Industrial participantes do Enade/2017, por Categoria Administrativa e Modalidade de ensino, de acordo com os conceitos por eles alcançados, segundo as Grandes Regiões brasileiras.
Dos 62 cursos participantes, 53 ( 85,5 \% ) eram ministrados em Instituições Privadas e nove ( 14,5 \% ), em Instituições Públicas.
Quanto à Modalidade de ensino, cinco cursos eram a Distância, e 57, Presenciais.
De acordo com as informações da Tabela 5.2, em termos nacionais, entre Instituições Públicas e Privadas, observa -se que, dos dois cursos avaliados com conceito 5, um era oferecido em IES Públicas, e outro, em IES Privadas.
Dos nove cursos participantes de IES Públicas, o conceito 4 foi o valor modal, atribuído a seis cursos ( 66,7 \% ).
Dos demais cursos de Instituições Públicas participantes, três cursos foram avaliados com conceitos 2, 3 e 5, um curso cada.
Nesta categoria, nenhum curso ficou Sem Conceito ( SC ) ou recebeu conceito 1.
Na rede Privada, o conceito modal foi 3, atribuído a 16 cursos, o correspondente a 30,2 \% dos 53 cursos da categoria.
Dos demais cursos participantes, dois ( 3,8 \% ) receberam conceito 1, e nove ( 17,0 \% ), conceito 2.
O conceito 4 foi atribuído a 14 cursos ( 26,4 \% ) e o conceito 5, a um curso ( 1,9 \% ).
Nesta Categoria Administrativa, 11 cursos ( 20,8 \% ) ficaram Sem Conceito ( SC ).
Dos 57 cursos participantes da modalidade Presencial, o conceito 4 foi o valor modal, atribuído a 20 cursos ( 35,1 \% ).
Dos demais cursos nesta Modalidade de ensino, 16 cursos ( 28,1 \% ) obtiveram conceito 3, oito cursos ( 14,0 \% ), conceito 2, dois cursos ( 3,5 \% ), conceito 5, e dois cursos ( 3,5 \% ) conceito 1.
Nesta modalidade, nove cursos ( 15,8 \% ) ficaram Sem Conceito ( SC ).
Dentre os cursos a Distância, o conceito modal foi 2, atribuído a dois cursos, o correspondente a 40,0 \% dos cinco cursos da modalidade.
Dos demais cursos participantes, um recebeu conceito 3 e outros dois ficaram Sem Conceito ( SC ).
Os demais conceitos não foram atribuídos a nenhum curso.
Na análise por região, observa -se que, na região Norte, as Instituições Privadas participaram com todos os quatro cursos da região, dos quais um ficou Sem Conceito ( SC ), um ficou com conceito 2, e dois ficaram com conceito 3, conceito modal.
Nenhum curso recebeu os demais conceitos, e todos os cursos oferecidos na região Norte eram presenciais.
Na região Nordeste, a rede privada também concentrou todos os quatro cursos participantes.
Dentre eles, três receberam conceito 2, e um ficou Sem Conceito ( SC ).
Todos os cursos na região eram presenciais.
Na região Sudeste, a proporção de cursos da rede privada, 73,5 \%, correspondendenderam a 25 dos 34 cursos participantes.
Nessa categoria e região, o conceito modal foi 3, atribuído a oito cursos.
Oito cursos ficaram Sem Conceito ( SC ), e nenhum recebeu conceito 5.
Os demais foram avaliados com conceito 1 ( dois cursos ), conceito 2 ( três cursos ) e conceito 4 ( quatro cursos ).
Dos nove cursos ( 26,5 \% ) oferecidos por Instituições Públicas ) na região Sudeste, o conceito modal foi 4, atribuído a seis cursos.
Os demais cursos receberam os conceitos 2, 3 e 5 ( um curso cada ).
Na região Sudeste, nesta categoria, nenhum curso ficou Sem Conceito ( SC ) ou recebeu conceito 1.
Na região Sudeste, a maioria dos cursos foram presenciais : 94,1 \%.
Para esses, a moda foi o conceito 4, com dez cursos.
Os demais cursos presenciais da região foram alocados aos conceitos 3 ( nove cursos ), 2 ( quatro cursos ), 1 ( dois cursos ) e 5 ( um curso ).
Seis cursos presenciais ficaram Sem Conceito ( SC ).
Os dois cursos a Distância da região Sudeste ficaram Sem Conceito ( SC ).
As Instituições Privadas concentraram todos os 19 cursos participantes da região Sul.
Desses, nove obtiveram conceito 4, o conceito modal.
Nesta combinação de Categoria Administrativa e Grande Região, os demais conceitos foram : 2 ( dois cursos ), 3 ( seis cursos ), e 5 ( um curso ).
Um curso ficou Sem Conceito ( SC ), e nenhum recebeu conceito 1.
Também na região Sul, a maioria dos cursos foram presenciais, correspondendo a 84,2 \%, em oposição aos 15,8 \% a Distância.
Para os cursos presenciais, a moda foi o conceito 4, com nove cursos.
Os demais cursos presenciais da região foram alocados aos conceitos 3 ( cinco cursos ) e 5 ( um curso ).
Nenhum curso Presencial foi alocado aos conceitos 1 e 2.
Somente um curso ficou Sem Conceito ( SC ).
Para os três cursos a Distância, a moda foi o conceito 2, com dois cursos ( 66,7 \% ).
O outro curso a Distância da região foi alocado ao conceito 3.
Nenhum curso a Distância da região foi alocado aos demais conceitos ou ficou Sem Conceito ( SC ).
Na região Centro-Oeste, houve apenas um curso participante, que era vinculado a Instituições Privadas e a ele foi atribuído o conceito 4.
Este mesmo curso era de Educação Presencial.
No Gráfico 5.2, apresenta -se a distribuição do Conceito Enade segundo a Categoria Administrativa da IES.
Os cursos em IES Públicas ( linha azul ) apresentam uma poligonal mais à direita do que os em IES Privadas ( linha verde ) e, consequentemente, uma distribuição de Conceitos Enade com valores maiores.
O que corrobora tal informação é que 22,2 \% dos cursos vinculados a IES Públicas receberam conceito 3 ou abaixo, ao passo que 64,3 \% dos 5.3 CONCEITOS POR ORGANIZAÇÃO ACADÊMICA E GRANDE REGIÃO Na Tabela 5.3, encontra -se a distribuição dos conceitos atribuídos aos cursos participantes do Enade/2017 na Área de Tecnologia em Gestão da Produção Industrial, por Organização Acadêmica, segundo as Grandes Regiões brasileiras.
Dos 62 cursos de Tecnologia em Gestão da Produção Industrial participantes, 19 eram oferecidos em Universidades, nove em Centros Universitários, 31 em Faculdades e três em CEFET/IFET.
Essa distribuição corresponde a, respectivamente, 30,6 \%, 14,5 \%, 50,0 \% e 4,8 \% dos cursos.
De acordo com os dados apresentados, os dois cursos avaliados com conceito 5 eram vinculados a Faculdades.
Os cursos vinculados a esse tipo de Organização Acadêmica tiveram o conceito 4 como conceito modal, atribuído a 13 cursos ( 41,9 \% ).
Os demais cursos vinculados a Faculdades receberam os conceitos 1 ( um curso ), 2 ( cinco cursos ), 3 ( oito cursos ) e conceito 5 ( dois cursos, como já mencionado ).
Ficaram Sem Conceito ( SC ) dois cursos.
Dentre os cursos em Universidades, os conceitos modais foram 3 e 4, atribuídos a cinco cursos cada ( 26,3 \% cada ).
Cinco cursos vinculados a esse tipo de Organização Acadêmica ficaram Sem Conceito ( SC ), e os outros cursos receberam os conceitos 1 ( um curso ) e 2 ( três cursos ).
Nenhum curso recebeu conceito 5.
Dos nove cursos mantidos por Centros Universitários, três ( 33,3 \% ) receberam conceito 3, o conceito modal, e dois ( 22,2 \% ) receberam conceito 2.
Quatro cursos ficaram Sem Conceito ( SC ), e nenhum recebeu conceito 1, 4 ou 5.
Dos três cursos mantidos por CEFET/IFET, dois ( 66,7 \% ) receberam conceito 4, conceito modal, e um ( 33,3 \% ) recebeu conceito 3.
Nenhum curso ficou Sem Conceito ( SC ) ou recebeu os conceitos 1, 2 ou 5.
Considerando -se separadamente as regiões brasileiras, verifica -se que, na região Norte, as Faculdades participaram com três dos quatro cursos participantes.
Dois cursos receberam conceito 3, conceito modal.
O outro curso recebeu o conceito 2.
Nenhum curso ficou Sem Conceito ( SC ) ou recebeu os demais conceitos.
Os Centros Universitários da região Norte foram representados por um curso, o qual ficou Sem Conceito ( SC ).
Não houve representação de cursos de Universidades e em CEFET/IFET na região Norte.
Na região Nordeste, as Universidades e os CEFET/IFET também não foram representados.
Dois cursos eram de Centros Universitários, dos quais um recebeu conceito 2 e outro ficou Sem Conceito ( SC ).
Os outros dois cursos da região Nordeste estavam vinculados a Faculdades e receberam conceito 2.
Na região Sudeste, as Faculdades concentraram 16 dos 34 cursos de Tecnologia em Gestão da Produção Industrial da região.
Dentre os cursos em Faculdades, nessa região, o conceito modal foi 4, atribuído a sete cursos.
Os demais cursos receberam os conceitos 1 ( um curso ), 2 ( dois cursos ), 3 ( quatro cursos ) e 5 ( um curso ).
Um curso ficou Sem Conceito ( SC ).
As Universidades participaram com 11 cursos na região Sudeste, dos quais quatro obtiveram os conceitos modais 2 e 3 ( dois cursos cada ).
Os conceitos 1 e 4 foram atribuídos a dois cursos, um curso cada.
Cinco cursos ficaram Sem Conceito ( SC ) nesta combinação de organização acadêmica e região e o conceito 5 não foi atribuído a nenhum curso.
Já os Centros Universitários foram representados por quatro cursos na região Sudeste, dos quais dois ficaram Sem Conceito ( SC ) e dois receberam conceito 3.
Os CEFET/IFET participaram com três, os quais receberam conceitos 4 ( dois cursos, o conceito modal ) e 3, ( um curso ).
Nenhum curso ficou Sem Conceito ( SC ), nem recebeu os demais conceitos.
Dos 19 cursos da região Sul, nove eram de Faculdades, para os quais o conceito modal foi 4, atribuído a cinco cursos.
Nesse tipo de organização, um curso ficou Sem Conceito ( SC ), e nenhum curso recebeu conceito 1 ou 2.
Os demais cursos receberam os conceitos 3 e 5 ( dois e um cursos, respectivamente ).
Dos oito cursos participantes de Universidades da região Sul, nenhum ficou Sem Conceito ( SC ) ou recebeu conceito 1 ou 5.
O conceito modal, 4, foi atribuído a quatro cursos.
Os conceitos restantes, 2 e 3, foram atribuídos, respectivamente, a um e a três cursos.
Quanto aos dois cursos vinculados a Centros Universitários na região Sul, um recebeu conceito 2, e outro, o conceito 3.
Nessa região, não foram avaliados cursos em CEFET/IFET.
Na região Centro-Oeste, o único curso representando recebeu conceito 4 e era vinculado a Faculdades.
No Gráfico 5.3, apresenta -se a distribuição cumulativa do Conceito Enade, segundo a Organização Acadêmica da IES.
Os cursos em CEFET/IFET ( linha roxa ) e em Faculdades ( linha vermelha ) apresentam uma poligonal mais à direita do que as demais, denotando uma melhor distribuição de Conceitos Enade.
No outro extremo, encontram -se os cursos oferecidos em Centros Universitários ( linha verde ) com a poligonal mais à esquerda e, consequentemente, uma distribuição de Conceitos Enade com menores valores.
IES de Universidades ( linha azul ) apresentam uma situação intermediária.
Nota -se que os cursos em IES de Centros Universitários ( linha verde ) receberam apenas os conceitos 2 e 3, pois a respectiva poligonal cumulativa parte de 0 \% do conceito 1 e alcança 100 \% no conceito 3.
Os de CEFET/IFET ( linha roxa ) apresentam a mesma configuração, mas os únicos conceitos atribuídos são 3 e 4.
Cursos em Universidades ( linha azul ) receberam conceito 4 ou valor inferior, pois a sua poligonal cumulativa alcança 100 \% no conceito 4.

% matched lemmas: 
\end{textsample}
